%======================================================================
\section{Theorem F: certified roots and continuation of the corrected equation}
\label{sec:certified}

Theorem~\CGref{thm:uniform_iteration} takes Newton steps with the fixed denominator $\lL$ and cutoffs that depend only on $T$. This section replaces the Newton steps by roots of the corrected equation that are certified index by index, and the fixed cutoffs by a cutoff that is raised index by index along a path on which the root cannot change branch. Everything is computed from $\tau_k$ and the primes.

\subsection{The sharp Euler phase}
\label{subsec:sharp_phase}
For $X\ge2$ let $A_{\le X}=\{n\le X:\ n\text{ a prime power}\}$ and
\begin{equation}
S_{\le X}=S_{A_{\le X}},\qquad F_{\le X}=\NO+S_{\le X},\qquad \sigma_{\le X}=\sum_{n\le X}|a_n|,\qquad M_{\le X}=\frac1\pi\sum_{n\le X}\frac{\Lambda(n)}{\sqrt n},
\label{eq:sharp_def}
\end{equation}
with $S_A$ and $a_n$ as in \CGeqref{eq:SA_def}. $S_{\le X}$ is a trigonometric polynomial, $|S_{\le X}|\le\sigma_{\le X}$ and $|S_{\le X}'|\le M_{\le X}$, since $a_n\log n=-\Lambda(n)/(\pi\sqrt n)$. It differs from $S_X=S_{A_X}$ by the powers $p^m>X$ of the primes $p\le X$, and
\begin{equation}
\sup_{t\in\R}|S_X(t)-S_{\le X}(t)|\le\frac1\pi\sum_{X^{1/2}<p\le X}\sum_{m\ge2}\frac1{mp^{m/2}}+\frac1\pi\sum_{p\le X^{1/2}}\sum_{m:\,p^m>X}\frac1{p^{m/2}}\ll1,
\label{eq:sharp_vs_full}
\end{equation}
by Theorem~\CGref{thm:known_mertens} for the first sum and $\sum_{p^m>X}p^{-m/2}\le X^{-1/2}(1-2^{-1/2})^{-1}$ for the second. A bounded change of the Euler phase changes the mean squares of Theorems~\CGref{thm:main} and \CGref{thm:computable_clock} by $O(1+\sqrt{\log(1/c)})$, so both theorems hold for $S_{\le X}$ with the same error terms.

\subsection{Monotonicity and certified roots}
\begin{lemma}[Certified monotonicity]\label{lem:monotone_sharp}
Let $0<\mathsf c<1$, $t_0>2\pi$ and $X\ge2$ with $M_{\le X}\le(1-\mathsf c)\lL(t_0)$. Then $F_{\le X}'(t)\ge\mathsf c\,\lL(t)$ for all $t\ge t_0$; in particular $F_{\le X}$ is strictly increasing on $[t_0,\infty)$ and each equation $F_{\le X}(t)=k-\frac12$ has at most one root there. The bound $|S'_{\le X}|\le M_{\le X}$ cannot be improved uniformly in $t$: $\sup_{t\in\R}\bigl(-S'_{\le X}(t)\bigr)=M_{\le X}$.
\end{lemma}

\begin{proof}
$F_{\le X}'=\lL+S'_{\le X}\ge\lL(t)-M_{\le X}\ge\lL(t)-(1-\mathsf c)\lL(t_0)\ge\mathsf c\lL(t)$ for $t\ge t_0$, because $\lL$ is increasing. For the last statement, $-S'_{\le X}(t)=\frac1\pi\sum_{n\le X}\Lambda(n)n^{-1/2}\cos(t\log n)$, and $\cos(t\log n)=\cos(mt\log p)$ for $n=p^m$. The numbers $\log p$, $p\le X$, are linearly independent over $\mathbb Q$ by unique factorization, so by Kronecker's theorem \cite[Ch.~XXIII]{HardyWright2008} there are real $t$ with $t\log p$ arbitrarily close to $0$ modulo $2\pi$ for all $p\le X$ simultaneously; at such $t$ every cosine is arbitrarily close to $1$.
\end{proof}

By the prime number theorem $\sum_{n\le X}\Lambda(n)n^{-1/2}\sim2\sqrt X$, so the largest admissible cutoff of Lemma~\CGref{lem:monotone_sharp} is
\begin{equation}
X_{\rm mon}(t_0)=\max\bigl\{X\ge1:\ M_{\le X}\le(1-\mathsf c)\lL(t_0)\bigr\}\sim\Bigl(\frac{(1-\mathsf c)\log(t_0/2\pi)}4\Bigr)^2 ,
\label{eq:Xmon}
\end{equation}
the scale $(\log T)^2$ of Proposition~\CGref{prop:envelope_equivalence}. Proposition~\CGref{prop:monotone} shows that this scale gains nothing in order; it is used here only as the starting point of a continuation.

\begin{lemma}[Certified root]\label{lem:certified_root}
Let $\mathscr I=[\mathsf a,\mathsf b]$, $F\in C^1(\mathscr I)$ with $F'\ge\mathsf d>0$ on $\mathscr I$, and $F(\mathsf a)<\mathsf v<F(\mathsf b)$. Then $F=\mathsf v$ has exactly one root $x^*\in\mathscr I$; for every $x\in\mathscr I$, $|x-x^*|\le|F(x)-\mathsf v|/\mathsf d$; and $j$ bisection steps locate $x^*$ in an interval of length $2^{-j}(\mathsf b-\mathsf a)$.
\end{lemma}

\begin{proof}
$F$ is strictly increasing and continuous, so the root exists and is unique, and the bisection invariant $F(\text{left})<\mathsf v<F(\text{right})$ is preserved. By the mean value theorem $|F(x)-\mathsf v|=|F(x)-F(x^*)|\ge\mathsf d|x-x^*|$.
\end{proof}

A root found by bisection has no residual: $F(x^*)=\mathsf v$ up to the evaluation error. A Newton step with the fixed denominator $\lL$, as in \CGeqref{eq:newton_iteration}, leaves the residual of the equation
\begin{equation}
\mathsf E_k(x)=F_{\le X}(x)-\Bigl(k-\frac12\Bigr),
\label{eq:equation_residual}
\end{equation}
computable from the clock and the primes. The following identity separates it from the defect at the ordinate.

\begin{proposition}[Residual identity]\label{prop:residual_identity}
For every finite set $A$ of prime powers, every $k$ and every real $x>2\pi$, with $u_k(A)=s_k-S_A(\gamma_k)$ and $\mathsf E_k(x)=\NO(x)+S_A(x)-(k-\frac12)$,
\begin{equation}
\NO(\gamma_k)-\NO(x)=\bigl(S_A(x)-S_A(\gamma_k)\bigr)-u_k(A)-\mathsf E_k(x).
\label{eq:residual_identity}
\end{equation}
If $\mathsf E_k(x)=0$ and $F_A=\NO+S_A$ satisfies $F_A'\ge\mathsf d>0$ on the segment between $x$ and $\gamma_k$, then $|x-\gamma_k|\le|u_k(A)|/\mathsf d$.
\end{proposition}

\begin{proof}
Insert the definitions: the right side is $S_A(x)-S_A(\gamma_k)-k+\frac12+\NO(\gamma_k)+S_A(\gamma_k)-\NO(x)-S_A(x)+k-\frac12$. If $\mathsf E_k(x)=0$, then $F_A(x)=k-\frac12$ and $F_A(\gamma_k)=k-\frac12-u_k(A)$, so $|u_k(A)|=|F_A(x)-F_A(\gamma_k)|\ge\mathsf d|x-\gamma_k|$.
\end{proof}

The three terms of \CGeqref{eq:residual_identity} are the transport, the defect of the primes not used, and the residual of the equation. Theorem~\CGref{thm:main} controls the second; the third vanishes for certified roots; the first is the term of Remark~\CGref{rem:newton}.

\subsection{Continuation without change of branch}
\begin{theorem}[Theorem F: continuation]\label{thm:continuation}
Let $F_0,F_1\in C^1(\R)$, $\mathsf v\in\R$, $x_0,x_1$ real with $F_0(x_0)=\mathsf v=F_1(x_1)$, let $\mathscr J$ be the closed segment between $x_0$ and $x_1$, and suppose $F_0'\ge\mathsf d$ and $F_1'\ge\mathsf d$ on $\mathscr J$ for some $\mathsf d>0$. For $\mathsf s\in[0,1]$ let $F_{\mathsf s}=(1-\mathsf s)F_0+\mathsf sF_1$. Then for every $\mathsf s$ the equation $F_{\mathsf s}=\mathsf v$ has exactly one root $x(\mathsf s)\in\mathscr J$; $x(0)=x_0$, $x(1)=x_1$; $\mathsf s\mapsto x(\mathsf s)$ is continuous on $[0,1]$ and $C^1$ where $x(\mathsf s)$ is interior to $\mathscr J$, with
\begin{equation}
\frac{dx}{d\mathsf s}=-\frac{F_1(x(\mathsf s))-F_0(x(\mathsf s))}{F_{\mathsf s}'(x(\mathsf s))},\qquad |x_1-x_0|\le\frac{\sup_{\mathscr J}|F_1-F_0|}{\mathsf d}.
\label{eq:continuation}
\end{equation}
\end{theorem}

\begin{proof}
If $x_0=x_1$ there is nothing to prove. Let $x_0<x_1$ (the other case is symmetric). On $\mathscr J$, $F_{\mathsf s}'\ge\mathsf d$ by convexity. Since $F_1$ is increasing on $\mathscr J$ and $F_1(x_1)=\mathsf v$, $F_1(x_0)\le\mathsf v$; similarly $F_0(x_1)\ge\mathsf v$. Hence $F_{\mathsf s}(x_0)=(1-\mathsf s)\mathsf v+\mathsf sF_1(x_0)\le\mathsf v$ and $F_{\mathsf s}(x_1)\ge\mathsf v$, so by monotonicity $F_{\mathsf s}=\mathsf v$ has exactly one root in $\mathscr J$. Continuity of $x(\mathsf s)$ follows from uniqueness and compactness, and the implicit function theorem gives the derivative at interior points. The bound is Lemma~\CGref{lem:certified_root} applied to $F_1$ at the point $x_0$: $|F_1(x_0)-\mathsf v|=|F_1(x_0)-F_0(x_0)|$.
\end{proof}

For the corrected equation, $F_0=F_{\le X}$, $F_1=F_{\le Y}$ with $Y>X$, $\mathsf v=k-\frac12$, and
\begin{equation}
\sup_{\R}|F_{\le Y}-F_{\le X}|\le\sum_{X<n\le Y}|a_n|.
\label{eq:continuation_step}
\end{equation}
Theorem~\CGref{thm:continuation} is applied along a ladder $X=Y_0<Y_1<\dots<Y_J$: at level $j$ the new root $x_j$ of $F_{\le Y_j}=k-\frac12$ is computed from $x_{j-1}$ and accepted only if $F_{\le Y_{j-1}}'$ and $F_{\le Y_j}'$ are bounded below by $\mathsf d_j>0$ on the segment between $x_{j-1}$ and $x_j$. Then $x_J$ is joined to the certified root of Lemma~\CGref{lem:monotone_sharp} by a continuous path of roots of the homotopies, with no change of branch and no crossing of a neighbouring root. If the condition fails at some level, the index keeps the last accepted root: the failure costs the gain of the higher cutoff but cannot create a wrong root. The condition is local, on a segment of length $\le\sum_{Y_{j-1}<n\le Y_j}|a_n|/\mathsf d_j$, and concerns finite trigonometric polynomials, so it can be certified by interval arithmetic or by a bound for $F''$. It is not limited by $X\le T^{1/2}$.

\begin{remark}[The Nyquist cutoff]\label{rem:nyquist}
Sampled at the clock points, $\sin(t\log n)$ advances by $\log n/\lL(\tau_k)+O(1/\tau_k)=2\pi\alpha_n(\tau_k)+O(1/\tau_k)$ per index, with
\begin{equation}
\alpha_n(t)=\frac{\log n}{\log(t/2\pi)},
\label{eq:ordinal_frequency}
\end{equation}
the ordinal frequency of the term $n$. The frequencies $\alpha$ and $1-\alpha$ give the same samples up to sign, so $\alpha_n=\frac12$, that is $n=(\tau/2\pi)^{1/2}$, is the Nyquist frequency of the clock. This is the scale at which the cross term of the primes switches on (Section~\CGref{subsec:onset}) and the limit $X\le T^{1/2}$ of Theorem~\CGref{thm:main}. A residual dominated by the prime powers just above a cutoff $X=(\tau/2\pi)^{c}$ has ordinal frequencies just above $c$, which fold to just below $1-c$; for $c=\frac12$ the residual alternates in sign from one index to the next.
\end{remark}

\subsection{A per-index cutoff rule}
\label{subsec:rule}
The rule used in Section~\CGref{subsec:rule_data} is explicit and uses no ordinate.
\begin{enumerate}[leftmargin=2em,label=(\roman*)]
\item \emph{Stage A.} $X^{\rm A}_k=X_{\rm mon}(\tau_k-1)$ with $\mathsf c=\frac1{10}$; the root $x^{\rm A}_k$ of $F_{\le X^{\rm A}_k}=k-\frac12$ is located by bisection (Lemmas~\CGref{lem:monotone_sharp} and \CGref{lem:certified_root}).
\item \emph{Rule R1.} Continuation from $x^{\rm A}_k$ along six geometric levels from $X^{\rm A}_k$ to $(\tau_k/2\pi)^{1/2}$, the Nyquist cutoff.
\item \emph{Rule R2.} The same, up to $(\tau_k/2\pi)^{c_{\rm G}(k)}$, with $c_{\rm G}$ given by Heuristic~\CGref{heur:gaussian_cutoff}.
\end{enumerate}
At each level the new root is computed by Newton's method with the true derivative and the acceptance condition is $F'>\frac1{10}\lL$ for the old and the new phase.

\begin{heuristic}[Gaussian cutoff]\label{heur:gaussian_cutoff}
Suppose that after Stage A the error in units of zeros is Gaussian with standard deviation $\hat\sigma$, independent of the phases. The term $n$ of the Euler phase, evaluated at the point instead of the ordinate, has a phase error $\log n\cdot\hat\sigma/\lL=2\pi\hat\sigma\alpha_n$, and including it reduces its contribution to the mean square if and only if $\mathbb E\cos(2\pi\hat\sigma\alpha_nZ_{\rm G})=e^{-2\pi^2\hat\sigma^2\alpha_n^2}>\frac12$, $Z_{\rm G}$ standard normal. This gives the cutoff exponent and the model value of $\hat\sigma$,
\begin{equation}
c_{\rm G}=\frac{\sqrt{2\log2}}{2\pi\hat\sigma},\qquad \hat\sigma_k^2=v_{\{n:\,X^{\rm A}_k<n\le\tau_k/2\pi\}}=\frac12\sum_{X^{\rm A}_k<n\le\tau_k/2\pi}a_n^2 ,
\label{eq:gaussian_cutoff}
\end{equation}
the second being the mean square of the omitted prime powers (Theorem~\CGref{thm:main}).
\end{heuristic}

The model overestimates the error after Stage A ($\hat\sigma=0.269$ against $0.199$ measured, Section~\CGref{subsec:rule_data}), so $c_{\rm G}$ is too small; Section~\CGref{subsec:beyond_nyquist} replaces $c_{\rm G}$ by fixed exponents $c\in[0.85,1.30]$.

%======================================================================
\section{Theorem G: phase coherence of the corrected points}
\label{sec:walk}

The mean square of Theorem~\CGref{thm:uniform_iteration} does not control the coherent part of the error at a fixed frequency. The frequency $1$ is singled out by the phase walk of the points.

\begin{definition}\label{def:walk}
For real points $x_1,\dots,x_K$ let $\mathscr P_j(x)=\sum_{k\le j}e^{ix_k}$ and
\begin{equation}
\mathscr C_K(x)=\frac1K\sum_{j\le K}\mathscr P_j(x)=\sum_{k\le K}\mathsf w_k\,e^{ix_k},\qquad \mathsf w_k=\frac{K-k+1}K,
\label{eq:walk_def}
\end{equation}
the centroid of the walk. For the ordinates and for zero-free points $x_k$ put $e_k=x_k-\gamma_k$ (as for the Newton points, $e_{j,k}=x^{(j)}_k-\gamma_k$) and $\mathsf y_k=\tau_k-x_k$.
\end{definition}

\begin{proposition}[Angular and radial parts]\label{prop:angular_radial}
\begin{equation}
\mathscr C_K(x)-\mathscr C_K(\gamma)=i\mathfrak A+\mathfrak B,\qquad\mathfrak A=\sum_{k\le K}\mathsf w_ke^{i\gamma_k}\sine_k,\qquad\mathfrak B=\sum_{k\le K}\mathsf w_ke^{i\gamma_k}(\cose_k-1),
\label{eq:angular_radial}
\end{equation}
and, relative to the clock, with $|\mathfrak r_k|\le|\mathsf y_k|^3/6$,
\begin{equation}
\mathscr C_K(x)-\mathscr C_K(\tau)=\sum_{k\le K}\mathsf w_ke^{i\tau_k}\Bigl(-i\mathsf y_k-\frac{\mathsf y_k^2}2+\mathfrak r_k\Bigr).
\label{eq:walk_clock_expansion}
\end{equation}
\end{proposition}

\begin{proof}
$e^{ix_k}-e^{i\gamma_k}=e^{i\gamma_k}(e^{ie_k}-1)$ and $e^{ix_k}-e^{i\tau_k}=e^{i\tau_k}(e^{-i\mathsf y_k}-1)$; Taylor's formula with remainder for $e^{-iy}$.
\end{proof}

A weighted mean of $e^{i\gamma_k}(e^{ie_k}-1)$ of order $10^{-5}$ displaces the centroid by $O(10)$ over $2\cdot10^6$ ordinates, while the mean square of $e_k$ is not affected. The quadratic term of \CGeqref{eq:walk_clock_expansion} is explicit for the one-step clock.

\begin{proposition}[Products and quotients]\label{prop:product_quotient}
For a finite set $A$ of prime powers let $\mathscr H_A(t)=\sum_{n\in A}|a_n|e^{it\log n}$, so that $S_A=-\Im\mathscr H_A$. For the one-step correction $\mathsf y_k=S_A(\tau_k)/\lL(\tau_k)$,
\begin{equation}
-\frac{\mathsf y_k^2}2=\frac{\Re\bigl(\mathscr H_A(\tau_k)^2\bigr)}{4\lL(\tau_k)^2}-\frac{|\mathscr H_A(\tau_k)|^2}{4\lL(\tau_k)^2},
\label{eq:product_quotient}
\end{equation}
and $e^{it}\mathscr H_A(t)^2$, $e^{it}\overline{\mathscr H_A(t)}^2$ and $e^{it}|\mathscr H_A(t)|^2$ are sums of exponentials of frequencies $1+\log(n_1n_2)$, $1-\log(n_1n_2)$ and
\begin{equation}
\upsilon_{n_1,n_2}=1+\log\frac{n_1}{n_2}\qquad(n_1,n_2\in A),
\label{eq:quotient_frequency}
\end{equation}
with coefficients $\pm\frac14|a_{n_1}a_{n_2}|$ after division by $4\lL^2$.
\end{proposition}

\begin{proof}
$(\Im z)^2=\frac12(|z|^2-\Re z^2)$ with $z=\mathscr H_A(\tau_k)$, and expansion of the products.
\end{proof}

The two kinds of terms behave differently along the clock.

\begin{proposition}[Quotient modes]\label{prop:quotient_modes}
Let $n_1\ne n_2$ be prime powers with $0<|\upsilon_{n_1,n_2}|\le\frac12\log(\tau_{k_1}/2\pi)$, and let $k_1<k_2$ with $\tau_{k_2}\le U$. Then
\begin{equation}
\Bigl|\sum_{k_1<k\le k_2}e^{i\upsilon_{n_1,n_2}\tau_k}\Bigr|\ll\frac{\log(U/2\pi)}{|\upsilon_{n_1,n_2}|},
\label{eq:quotient_mode}
\end{equation}
and for any nonincreasing weights $0\le\mathsf g_k\le\mathsf g_{k_1}$ the weighted sum is $\ll\mathsf g_{k_1}\log(U/2\pi)/|\upsilon_{n_1,n_2}|$. In particular the contribution of the pair to the quadratic part of \CGeqref{eq:walk_clock_expansion} with $A$ fixed,
\begin{equation}
\mathscr R_{n_1,n_2}=-\frac{|a_{n_1}a_{n_2}|}4\sum_{k\le K}\frac{\mathsf w_k}{\lL(\tau_k)^2}e^{i\upsilon_{n_1,n_2}\tau_k},
\label{eq:quotient_contribution}
\end{equation}
is bounded by $\ll|a_{n_1}a_{n_2}|\log(\tau_K/2\pi)/(|\upsilon_{n_1,n_2}|\lL(\tau_1)^2)$ uniformly in $K$, and it oscillates with period
\begin{equation}
P_{n_1,n_2}(T)=\frac{\log(T/2\pi)}{|\upsilon_{n_1,n_2}|}
\label{eq:quotient_period}
\end{equation}
indices at height $T$.
\end{proposition}

\begin{proof}
Let $\Omega(x)=\upsilon\tau_x/2\pi$ with $\upsilon=\upsilon_{n_1,n_2}$. By \CGeqref{eq:clock}, $\Omega'(x)=\upsilon/\log(\tau_x/2\pi)$ is continuous and monotone, and $0<|\Omega'|\le\frac12$ on $[k_1,k_2]$, so $\|\Omega'\|=|\Omega'|\ge|\upsilon|/\log(U/2\pi)$; Theorem~\CGref{thm:known_KL} gives \CGeqref{eq:quotient_mode}. The weighted sum follows by summation by parts, and $\mathsf w_k/\lL(\tau_k)^2$ is nonincreasing in $k$. The phase advances by $\upsilon/\lL(\tau)+O(1/\tau)$ per index, which gives the period.
\end{proof}

For a fixed set $A$ every quotient mode is bounded. If the cutoff grows with the height, pairs with $|\upsilon_{n_1,n_2}|$ arbitrarily small enter, and the bound degenerates: $|\upsilon_{n_1,n_2}|$ is small exactly when $n_2/n_1$ is close to $e$. Among the first sixteen convergents of the continued fraction $e=[2;1,2,1,1,4,1,1,6,\dots]$ \cite{Olds1970}, those with numerator and denominator both prime powers are
\begin{equation}
\frac83,\quad\frac{11}4,\quad\frac{19}7,\quad\frac{193}{71},\quad\frac{49171}{18089},\qquad |\upsilon|=1.9\cdot10^{-2},\ 1.2\cdot10^{-2},\ 1.5\cdot10^{-3},\ 1.0\cdot10^{-5},\ 1.0\cdot10^{-10}
\label{eq:e_convergents}
\end{equation}
(computed by \texttt{opc\_e\_convergents\_v15.py}). The pair $(71,193)$ has period $P_{71,193}\approx1.09\cdot10^6$ indices at $T=5\cdot10^5$, half the length of the table, and dominates the drift of the corrected points (Section~\CGref{subsec:walk_data}).

\begin{remark}[Product modes]\label{rem:product_modes}
The frequencies $1\pm\log(n_1n_2)$ of \CGeqref{eq:product_quotient} are stationary for the sum over the clock when $\log(n_1n_2)\mp1=\log(\tau/2\pi)$, that is $n_1n_2=e^{\pm1}\tau/2\pi$, the analogue of the resonance $pq\approx\tau/2\pi$ of Proposition~\CGref{prop:pair_resonance} shifted by the factor $e^{\pm1}$. A term $\mathsf b\,e^{-i\tau_k\log n}$ in the walk produces at the stationary index a jump of size $|\mathsf b|\lL(\tau)\sqrt{2\pi\tau}$ spread over $\asymp\lL(\tau)\sqrt{2\pi\tau}$ indices; this follows from Theorem~\CGref{thm:known_stationary} as in the evaluation of Section~\CGref{subsec:onset}, with the same heuristic status (endpoint terms and the variation of the amplitude neglected). The first-order term $-i\mathsf y_k$ has frequencies $1\pm\log n$, $n\in A$, which are stationary only for $n=e^{\pm1}\tau/2\pi$; for cutoffs $X\le(\tau/2\pi)^{1/2}$ these do not occur.
\end{remark}

%======================================================================
\section{Theorem H: separation of the ordinates and zeros on the critical line}
\label{sec:separation}

The statements of this section connect the zero-free constructions with the position of the zeros relative to the critical line. They are elementary, but they identify exactly what a construction must achieve. Let $Z(t)=e^{i\vartheta(t)}\zeta(\frac12+it)$ be Hardy's function, real for real $t$ \cite[\S4.17]{Titchmarsh1986}% VERIFICAR section
, whose zeros are the ordinates of the zeros on the critical line.

\begin{proposition}[Isolated ordinates]\label{prop:isolated}
If $\gamma_{k-1}<\gamma_k<\gamma_{k+1}$ in the list of ordinates counted with multiplicity, then $\rho_k$ is a simple zero on the critical line, and $Z$ changes sign at $\gamma_k$. In particular, all zeros with $0<\gamma\le t$ are simple and lie on the critical line if and only if $\gamma_1<\gamma_2<\dots<\gamma_{N(t)}$.
\end{proposition}

\begin{proof}
By the functional equation and $\zeta(\bar s)=\overline{\zeta(s)}$, $1-\bar\rho$ is a zero with the multiplicity of $\rho$ \cite[\S2.1]{Titchmarsh1986}. If $\beta_k\ne\frac12$, the zero $1-\bar\rho_k\ne\rho_k$ has ordinate $\gamma_k$; if $\rho_k$ is multiple, it is counted at least twice. In both cases $\gamma_k$ occurs at least twice in the list. A simple zero on the critical line is a simple zero of the real analytic function $Z$, hence a sign change. The converse is clear.
\end{proof}

\begin{theorem}[Theorem H: local separation]\label{thm:local_separation}
Let $x_{k-1},x_k,x_{k+1}$ be real, $e_j=x_j-\gamma_j$, $\bar x_j=\frac12(x_j+x_{j+1})$, and $\mathsf h>0$. If
\begin{equation}
|e_j|<\mathsf h\ (j=k-1,k,k+1),\qquad x_k-x_{k-1}>2\mathsf h,\qquad x_{k+1}-x_k>2\mathsf h,
\label{eq:local_separation}
\end{equation}
then $\gamma_k$ is the only element of the list of ordinates in $(\bar x_{k-1},\bar x_k)$, $\rho_k$ is a simple zero on the critical line, and $Z(\bar x_{k-1})Z(\bar x_k)<0$.
\end{theorem}

\begin{proof}
$\gamma_k>x_k-\mathsf h>x_k-\frac12(x_k-x_{k-1})=\bar x_{k-1}$, and similarly $\gamma_k<\bar x_k$, $\gamma_{k-1}<\bar x_{k-1}$, $\gamma_{k+1}>\bar x_k$. The list is nondecreasing, so no other element lies in $(\bar x_{k-1},\bar x_k)$, and $\gamma_{k-1}<\gamma_k<\gamma_{k+1}$. Proposition~\CGref{prop:isolated} applies; $Z$ has exactly one zero in the interval, of odd order, and does not vanish at its endpoints.
\end{proof}

\begin{corollary}[Proportion of simple zeros on the line]\label{cor:proportion}
Let $x_k$ be any real points, $I$ a set of consecutive indices, $I'$ the set $I$ enlarged by one index at each end, and $\mathsf h>0$. The number of simple zeros on the critical line with ordinate in $[\min_I\bar x_{k-1},\max_I\bar x_k]$ is at least
\begin{equation}
\#I-\frac3{\mathsf h^2}\sum_{j\in I'}e_j^2-2\,\#\{j\in I':\ x_{j+1}-x_j\le2\mathsf h\}.
\label{eq:proportion}
\end{equation}
\end{corollary}

\begin{proof}
An index $k\in I$ fails \CGeqref{eq:local_separation} only if $|e_j|\ge\mathsf h$ for some $j\in\{k-1,k,k+1\}$, or one of the two gaps at $k$ is $\le2\mathsf h$; each $j$ serves at most three indices $k$ and each gap at most two. Chebyshev's inequality bounds $\#\{j:|e_j|\ge\mathsf h\}$ by $\mathsf h^{-2}\sume_j^2$. The intervals $(\bar x_{k-1},\bar x_k)$ for distinct $k$ are disjoint, and each contains a distinct simple zero on the line by Theorem~\CGref{thm:local_separation}.
\end{proof}

With the clock corrected $m^*$ times (Theorem~\CGref{thm:uniform_iteration}) and $\mathsf h=\mathsf u/\lL(T)$, \CGeqref{eq:proportion} gives at least
\begin{equation}
\Bigl(1-\frac{3C'}{\mathsf u^2}(1+o(1))-2\,\frac{\#\{k\in\IT:\ \lL(T)(x^{(m^*)}_{k+1}-x^{(m^*)}_k)\le2\mathsf u\}}{\#\IT}\Bigr)\#\IT
\label{eq:proportion_iteration}
\end{equation}
simple zeros on the line with ordinates near $(T,2T]$, where $C'$ is the constant of \CGeqref{eq:final_clock}. The last term concerns only the zero-free points. A positive proportion follows from this route as soon as $C'$ is explicit and small compared with $\mathsf u^2$ and the gaps of the corrected clock are controlled; neither is proved here. The proportion $\frac5{12}$ of \cite{PrattRoblesZaharescuZeindler2020} (Selberg proved a positive proportion \cite{Selberg1942}) is exceeded if the bracket in \CGeqref{eq:proportion_iteration} exceeds $\frac5{12}$.

\begin{proposition}[Verification with the corrected points]\label{prop:verification}
Let $x_k$ be zero-free points with $\bar x_{k_1}<\bar x_{k_1+1}<\dots<\bar x_{k_2}$, and suppose $Z(\bar x_{k-1})Z(\bar x_k)<0$ for $k_1<k\le k_2$ (the \emph{corrected Gram law} at these indices). Then there are at least $k_2-k_1$ zeros on the critical line with ordinate in $(\bar x_{k_1},\bar x_{k_2})$. If moreover $N(\bar x_{k_2})-N(\bar x_{k_1})=k_2-k_1$, then all zeros with ordinate in this interval are simple and lie on the critical line.
\end{proposition}

\begin{proof}
Each interval $(\bar x_{k-1},\bar x_k)$ contains a sign change of $Z$, hence a zero on the line of odd order; the intervals are disjoint. If there are exactly $k_2-k_1$ zeros counted with multiplicity, each of them is one of these, of order one.
\end{proof}

For the bare clock the midpoints $\bar\tau_k$ lie within $O(\tau_k^{-1})$ of the Gram points, by \CGeqref{eq:clock}, and the corrected Gram law at $x=\tau$ is Gram's law \cite{Gram1903}. The count $N(\bar x_{k_2})$ is computable by Turing's method \cite{Turing1953}, so Proposition~\CGref{prop:verification} is a verification procedure in which the corrected points replace the Gram points; its exceptions are the indices at which the law fails (Section~\CGref{subsec:gram_data}). For all ordinates up to $3\cdot10^{12}$ the hypothesis of Proposition~\CGref{prop:isolated} holds \cite{PlattTrudgian2021}, which is used in Section~\CGref{subsec:gram_data} to evaluate the law from the ordinates.

